\documentstyle[12pt,fullpage]{article}
\setcounter{page}{1}

\begin{document}

\title{EE 150 Quiz \#4}
\author{Due: Nov 1, 1996 (in class)}
\date{}
\maketitle

Given two arrays defined int array1[SIZE] and int array2[SIZE].  We
know that if $i$ represents the {\em index} of a particular element
in array1, then array1[i] represents the element itself.

Write a function {\bf arrayMultiply} that takes two arrays (m1 and
m2) as parameters, and, for all values in the arrays, sets the value 
of each position $i$ in the m1 array  to equal the {\em product} of
the original value in m1[i] and the value of m2[i].   The function
should look like this ....

\begin{verbatim}
int arrayMultiply( int m1[], int m2[], SIZE )
{
 .... fill in here ... 
}

For example, if SIZE were 5, and the values stored in array1 were 0 1
2 3 4, and the values stored in array2 where 9 8 7 6 5, then after a
function call arrayMultiply( array1, array2, 5 ), array1 would contain
0, 8, 14, 18, 20. 
\end{verbatim}

\end{document}
